% =====================================================================
\chapter{Espacios de probabilidad}
\label{tema:3}
% =====================================================================

Este tema da la definición axiomática de probabilidad, deduce de los axiomas las reglas de cálculo y estudia el caso más importante para empezar, el modelo de Laplace, en el que calcular una probabilidad se reduce a contar. Se siguen la sección §1.2 de Bertsekas y Tsitsiklis~\cite{bt} (clase~1 de MIT 6.041~\cite{mit6041f10}), el capítulo~2 de Ross~\cite{ross} (§2.2 a §2.5), las secciones §1.2, §1.3 y §1.6 de Blitzstein y Hwang~\cite{bh} y el capítulo~16, «Events and Probability Spaces», de~\cite{mcs}, que incluye el problema del cumpleaños (§16.4).

\section{Experimentos, espacio muestral y sucesos}

\begin{definicion}[Espacio muestral y sucesos]
Un \emph{experimento aleatorio} es un procedimiento cuyo resultado no puede predecirse con certeza. El \emph{espacio muestral} $\Omega$ es el conjunto de todos sus resultados posibles. Un \emph{suceso} es un subconjunto $A\subseteq\Omega$; se dice que $A$ \emph{ocurre} si el resultado obtenido pertenece a $A$. Los sucesos $\{\omega\}$ se llaman \emph{elementales}.
\end{definicion}

Las operaciones con conjuntos del módulo~0.2 se traducen en operaciones con sucesos: $A\cup B$ es «ocurre $A$ o $B$»; $A\cap B$, «ocurren $A$ y $B$»; $A^{c}=\Omega\setminus A$, «no ocurre $A$»; $A\subseteq B$, «si ocurre $A$, ocurre $B$». Dos sucesos son \emph{incompatibles} (o mutuamente excluyentes) si $A\cap B=\emptyset$. Las leyes de De Morgan dan, por ejemplo, $(A\cup B)^{c}=A^{c}\cap B^{c}$: «no ocurre ninguno de los dos».

\begin{ejemplo}\label{ej:t3-omega}
Al lanzar dos dados distinguibles, $\Omega=\{(i,j):i,j\in[6]\}$ tiene 36 resultados. El suceso «la suma es 7» es $\{(1,6),(2,5),(3,4),(4,3),(5,2),(6,1)\}$. Un modelo de lenguaje que genera dos tokens de un vocabulario $\V$ tiene espacio muestral $\V^{2}$, y «el primer token es \texttt{el}» es el suceso $\{\texttt{el}\}\times\V$.
\end{ejemplo}

\section{Axiomas de Kolmogorov}

\begin{definicion}[Espacio de probabilidad]\label{def:t3-axiomas}
Una \emph{ley de probabilidad} sobre $\Omega$ es una función $\PP$ que asigna a cada suceso $A$ un número real $\PP(A)$ y que cumple:
\begin{enumerate}
  \item \emph{no negatividad}: $\PP(A)\ge0$ para todo suceso $A$;
  \item \emph{normalización}: $\PP(\Omega)=1$;
  \item \emph{aditividad numerable}: si $A_1,A_2,\dots$ son sucesos incompatibles dos a dos, $\PP\bigl(\bigcup_{i\ge1}A_i\bigr)=\sum_{i\ge1}\PP(A_i)$.
\end{enumerate}
El par $(\Omega,\PP)$ se llama \emph{espacio de probabilidad}.
\end{definicion}

\begin{observacion}
Si $\Omega$ es finito o numerable, todo subconjunto puede ser un suceso. Si $\Omega$ es no numerable, como un intervalo de $\R$, no es posible en general asignar probabilidad a todos los subconjuntos de forma coherente, y los sucesos se restringen a una familia adecuada (una $\sigma$-álgebra). En este módulo esa sutileza no interviene: todos los sucesos que aparecen (intervalos, uniones de intervalos, regiones del plano) son admisibles~\cite[§1.6]{bh}.
\end{observacion}

\section{Propiedades que se deducen de los axiomas}

\begin{proposicion}\label{prop:t3-props}
Sean $A$, $B$ y $A_1,\dots,A_n$ sucesos. Entonces:
\begin{enumerate}
  \item $\PP(\emptyset)=0$;
  \item \emph{aditividad finita}: si $A_1,\dots,A_n$ son incompatibles dos a dos, $\PP\bigl(\bigcup_{i=1}^{n}A_i\bigr)=\sum_{i=1}^{n}\PP(A_i)$;
  \item $\PP(A^{c})=1-\PP(A)$;
  \item si $A\subseteq B$, entonces $\PP(B\setminus A)=\PP(B)-\PP(A)$ y $\PP(A)\le\PP(B)$;
  \item $0\le\PP(A)\le1$;
  \item $\PP(A\cup B)=\PP(A)+\PP(B)-\PP(A\cap B)$.
\end{enumerate}
\end{proposicion}

\begin{proof}
(i) Los sucesos $A_1=\Omega$, $A_i=\emptyset$ ($i\ge2$) son incompatibles dos a dos y su unión es $\Omega$. Por el axioma (iii), $1=\PP(\Omega)=1+\sum_{i\ge2}\PP(\emptyset)$, luego $\sum_{i\ge2}\PP(\emptyset)=0$; como los sumandos son iguales y no negativos, $\PP(\emptyset)=0$.

(ii) Se completa la lista con $A_i=\emptyset$ para $i>n$ y se aplican el axioma (iii) y el apartado~(i).

(iii) $A$ y $A^{c}$ son incompatibles y su unión es $\Omega$: $\PP(A)+\PP(A^{c})=1$.

(iv) Si $A\subseteq B$, $B$ es la unión disjunta de $A$ y $B\setminus A$; por (ii), $\PP(B)=\PP(A)+\PP(B\setminus A)$, y como $\PP(B\setminus A)\ge0$, $\PP(A)\le\PP(B)$.

(v) Por (iv) con $A\subseteq\Omega$, $\PP(A)\le\PP(\Omega)=1$.

(vi) $A\cup B$ es la unión disjunta de $A$ y $B\setminus(A\cap B)$, y $A\cap B\subseteq B$. Por (ii) y (iv), $\PP(A\cup B)=\PP(A)+\PP(B)-\PP(A\cap B)$.
\end{proof}

\begin{teorema}[Desigualdad de Boole o cota de la unión]\label{teo:t3-boole}
Para sucesos cualesquiera $A_1,\dots,A_n$,
\(
\PP\bigl(\bigcup_{i=1}^{n}A_i\bigr)\le\sum_{i=1}^{n}\PP(A_i).
\)
\end{teorema}

\begin{proof}
Por inducción sobre $n$. El caso $n=1$ es trivial. Si es cierto para $n-1$, por la \cref{prop:t3-props}(vi) aplicada a $B=A_1\cup\dots\cup A_{n-1}$ y $A_n$,
\[
\PP(B\cup A_n)=\PP(B)+\PP(A_n)-\PP(B\cap A_n)\le\PP(B)+\PP(A_n)\le\sum_{i=1}^{n-1}\PP(A_i)+\PP(A_n).\qedhere
\]
\end{proof}

La fórmula de inclusión--exclusión del \cref{teo:t2-ie} tiene una versión probabilística con la misma estructura; para tres sucesos,
\[
\PP(A\cup B\cup C)=\PP(A)+\PP(B)+\PP(C)-\PP(A\cap B)-\PP(A\cap C)-\PP(B\cap C)+\PP(A\cap B\cap C),
\]
y se demuestra a partir de (vi) exactamente como el \cref{teo:t1-ie}.

\section{Espacios discretos y modelo de Laplace}

\begin{proposicion}[Espacios discretos]\label{prop:t3-discreto}
Sea $\Omega$ finito o numerable y $p\colon\Omega\to[0,1]$ con $\sum_{\omega\in\Omega}p(\omega)=1$. Entonces $\PP(A)=\sum_{\omega\in A}p(\omega)$ define una ley de probabilidad. Recíprocamente, toda ley de probabilidad sobre $\Omega$ es de esta forma, con $p(\omega)=\PP(\{\omega\})$.
\end{proposicion}

\begin{proof}
La no negatividad y la normalización son inmediatas. Si $A_1,A_2,\dots$ son disjuntos, la suma de $p$ sobre $\bigcup_iA_i$ puede reordenarse agrupando por conjuntos, porque es una serie de términos no negativos, y da $\sum_i\PP(A_i)$. Para el recíproco, todo $A\subseteq\Omega$ es la unión disjunta, finita o numerable, de los sucesos elementales $\{\omega\}$ con $\omega\in A$; por el axioma~(iii), $\PP(A)=\sum_{\omega\in A}\PP(\{\omega\})$.
\end{proof}

\begin{definicion}[Modelo de Laplace]\label{def:t3-laplace}
Si $\Omega$ es finito y todos los resultados son equiprobables, $p(\omega)=1/\abs{\Omega}$, y entonces
\[
\PP(A)=\frac{\abs{A}}{\abs{\Omega}}=\frac{\text{casos favorables}}{\text{casos posibles}}.
\]
\end{definicion}

El modelo de Laplace es una \emph{hipótesis} sobre el experimento, no una consecuencia de los axiomas: es razonable para dados, monedas y cartas bien mezcladas, pero no, por ejemplo, para la suma de dos dados, cuyos 11 valores no son equiprobables. La técnica habitual consiste en elegir un espacio muestral en el que la equiprobabilidad sea evidente, normalmente uno de resultados \emph{ordenados} y distinguibles, aunque los objetos del problema parezcan indistinguibles.

\begin{ejemplo}[Ordenado o no ordenado]\label{ej:t3-ases}
Se extraen dos cartas de una baraja francesa sin reemplazamiento. Con resultados ordenados, $\abs{\Omega}=52\cdot51$ y los pares de ases son $4\cdot3$, luego $\PP(\text{dos ases})=\frac{12}{2652}=\frac{1}{221}$. Con resultados no ordenados, $\abs{\Omega}=\binom{52}{2}=1326$ y los favorables son $\binom{4}{2}=6$: de nuevo $\frac{6}{1326}=\frac1{221}$. Ambos modelos son válidos porque cada subconjunto corresponde exactamente a dos pares ordenados; lo que no es válido es mezclar recuentos ordenados en el numerador con no ordenados en el denominador.
\end{ejemplo}

\section{El problema del cumpleaños}

\begin{proposicion}\label{prop:t3-cumple}
Se asigna a cada una de $n$ personas un día entre $N$, de forma uniforme e independiente (todas las $N^{n}$ asignaciones son equiprobables). La probabilidad de que dos o más personas compartan día es
\[
p_n=1-\frac{N^{\underline{n}}}{N^{n}}=1-\prod_{i=0}^{n-1}\Bigl(1-\frac{i}{N}\Bigr),
\]
y además $1-p_n\le\exp\!\bigl(-\tfrac{n(n-1)}{2N}\bigr)$.
\end{proposicion}

\begin{proof}
El suceso complementario, «todos los días son distintos», está formado por las asignaciones inyectivas, que son $N^{\underline{n}}$ (\cref{teo:t1-variaciones}); se aplica el modelo de Laplace y la regla del complemento. Para la cota se usa la desigualdad $1-x\le e^{-x}$, válida para todo $x$ real (la función $g(x)=e^{-x}-1+x$ tiene $g'(x)=1-e^{-x}$, que se anula solo en $x=0$, donde $g$ alcanza su mínimo $g(0)=0$). Entonces
\[
1-p_n=\prod_{i=0}^{n-1}\Bigl(1-\frac iN\Bigr)\le\prod_{i=0}^{n-1}e^{-i/N}=\exp\Bigl(-\frac1N\sum_{i=0}^{n-1}i\Bigr)=\exp\Bigl(-\frac{n(n-1)}{2N}\Bigr).\qedhere
\]
\end{proof}

\begin{ejemplo}
Con $N=365$ y $n=23$, el producto da $1-p_{23}\approx0{,}4927$, luego $p_{23}\approx0{,}5073$: en un grupo de 23 personas es más probable que haya dos con el mismo cumpleaños que lo contrario. La cota da $1-p_{23}\le e^{-253/365}\approx0{,}49999$, suficiente para asegurar $p_{23}>1/2$. La razón de que baste un número tan pequeño es que lo que importa es el número de \emph{pares}, $\binom{23}{2}=253$, no el de personas~\cite[§16.4]{mcs}.
\end{ejemplo}

\begin{figure}[h]
\centering
\begin{tikzpicture}[x=0.11cm,y=3.2cm]
  \draw[->] (0,0) -- (63,0) node[right] {\small $n$};
  \draw[->] (0,0) -- (0,1.12) node[above] {\small $p_n$};
  \foreach \x in {10,20,30,40,50,60} \draw (\x,0.02) -- (\x,-0.02) node[below] {\scriptsize \x};
  \foreach \y in {0.5,1} \draw (0.6,\y) -- (-0.6,\y) node[left] {\scriptsize \y};
  \draw[dashed,gray] (0,0.5) -- (60,0.5);
  \draw[dashed,gray] (23,0) -- (23,0.5073);
  \draw[thick] plot[smooth] coordinates {(1,0) (5,0.0271) (10,0.1169) (15,0.2529) (20,0.4114) (23,0.5073) (25,0.5687) (30,0.7063) (35,0.8144) (40,0.8912) (45,0.9410) (50,0.9704) (55,0.9863) (60,0.9941)};
  \fill (23,0.5073) circle (0.9pt);
\end{tikzpicture}
\caption{Probabilidad $p_n$ de que al menos dos de $n$ personas compartan cumpleaños ($N=365$).}
\label{fig:t3-cumple}
\end{figure}

\begin{observacion}[Conexión con el curso]
El problema del cumpleaños es el modelo de las \emph{colisiones} de funciones \emph{hash}: si $n$ documentos se asignan uniformemente a $N$ valores, la probabilidad de colisión es $p_n$, y la cota $1-p_n\le e^{-n(n-1)/(2N)}$ muestra que las colisiones aparecen cuando $n$ es del orden de $\sqrt{N}$. Este razonamiento interviene en la deduplicación de corpus con MinHash del módulo~16. El espacio muestral $\V^{T}$ de secuencias de tokens y la regla del complemento son la base del módulo~1, y la cota de la unión reaparece cuando se acota la probabilidad de que falle alguno de muchos sucesos.
\end{observacion}

\begin{observacion}[Errores frecuentes]
(1) Suponer equiprobabilidad donde no la hay, como en las sumas de dos dados. (2) Mezclar recuentos ordenados y no ordenados en el mismo cociente (\cref{ej:t3-ases}). (3) Sumar probabilidades de sucesos que no son incompatibles: $\PP(A\cup B)=\PP(A)+\PP(B)$ solo si $A\cap B=\emptyset$; en general, la suma es solo una cota superior (\cref{teo:t3-boole}). (4) Tratar como indistinguibles objetos físicamente distintos (dos dados iguales siguen siendo dos dados). (5) Confundir «alguien comparte cumpleaños con usted» (\cref{eje:3-cumple}) con «dos personas cualesquiera comparten cumpleaños».
\end{observacion}

\begin{lecturas}
\item Bertsekas y Tsitsiklis~\cite{bt}, §1.1 y §1.2 (clase~1 de MIT 6.041~\cite{mit6041f10}, «Probability models and axioms»).
\item Ross~\cite{ross}, cap.~2 «Axioms of Probability», §2.2 a §2.5.
\item Blitzstein y Hwang~\cite{bh}, §1.2 «Sample spaces and Pebble World», §1.3 «Naive definition of probability» y §1.6 «Non-naive definition of probability».
\item Lehman, Leighton y Meyer~\cite{mcs}, cap.~16 «Events and Probability Spaces», §16.4 «The Birthday Principle» y §16.5 «Set Theory and Probability».
\item Grinstead y Snell~\cite{gs}, §1.2 «Discrete Probability Distributions».
\end{lecturas}

\section*{Ejercicios del tema 3}
\addcontentsline{toc}{section}{Ejercicios del tema 3}

\begin{ejercicio}{1}{eje:3-espacio}
Describa un espacio muestral adecuado y su cardinal para: (a) tres lanzamientos de una moneda; (b) una mano de 5 cartas; (c) la generación de dos tokens con un vocabulario de $K$ tokens; (d) el número de lanzamientos de una moneda hasta la primera cara. ¿En cuáles es razonable el modelo de Laplace?
\end{ejercicio}

\begin{ejercicio}{1}{eje:3-sucesos}
Sean $A$, $B$ y $C$ sucesos. Exprese con uniones, intersecciones y complementarios: (a) ocurre al menos uno; (b) ocurren los tres; (c) no ocurre ninguno; (d) ocurre exactamente uno; (e) ocurre como máximo uno; (f) ocurren $A$ y $B$, pero no $C$.
\fuente{adaptado de MIT 6.041~\cite{mit6041f10}, \emph{Problem Set}~1, problema~1.}
\end{ejercicio}

\begin{ejercicio}{2}{eje:3-monedas}
Se lanza tres veces una moneda equilibrada. Calcule la probabilidad de (a) obtener CCC; (b) obtener CXC; (c) obtener exactamente dos caras; (d) obtener al menos tantas caras como cruces.
\fuente{adaptado de MIT 6.041~\cite{mit6041f10}, \emph{Problem Set}~1, problema~2.}
\end{ejercicio}

\begin{ejercicio}{2}{eje:3-dados}
Se lanzan dos dados equilibrados. Calcule la probabilidad de que (a) la suma sea 7; (b) salga al menos un 6; (c) el máximo de los dos valores sea~4. En (b), compare con el resultado erróneo $\frac16+\frac16$ y explique el error.
\end{ejercicio}

\begin{ejercicio}{2}{eje:3-cuatro-dados}
Se lanzan cuatro dados equilibrados. Sin calcular ninguna de las dos probabilidades por separado, decida cuál es mayor, la de que la suma sea 21 o la de que sea 22. Después calcule ambas.
\fuente{adaptado de Harvard Stat 110~\cite{stat110}, \emph{Strategic Practice}~1, problema~1(a).}
\end{ejercicio}

\begin{ejercicio}{2}{eje:3-horarios}
Una facultad asigna cada asignatura, al azar y de forma independiente, a una de 10 franjas horarias. Un estudiante elige 3 asignaturas. ¿Qué probabilidad hay de que dos o más coincidan en la misma franja?
\fuente{adaptado de Harvard Stat 110~\cite{stat110}, \emph{Homework}~1, problema~3.}
\end{ejercicio}

\begin{ejercicio}{2}{eje:3-robos}
En una ciudad con 6 distritos se producen 6 robos; cada uno ocurre en un distrito elegido al azar, independientemente de los demás. ¿Qué probabilidad hay de que algún distrito haya sufrido más de un robo?
\fuente{adaptado de Harvard Stat 110~\cite{stat110}, \emph{Homework}~1, problema~4.}
\end{ejercicio}

\begin{ejercicio}{2}{eje:3-seis-cartas}
Se reparte una mano de 6 cartas de una baraja francesa. Se llama \emph{doble pareja} a una mano con exactamente dos valores repetidos dos veces y las otras dos cartas de valores distintos entre sí y de los anteriores, y \emph{trío} a una mano con un valor repetido tres veces y las otras tres cartas de valores distintos entre sí y del anterior. ¿Cuál es más probable? Calcule ambas probabilidades.
\fuente{adaptado de MIT 18.05~\cite{mit1805}, \emph{Problem Set}~1, problema~1.}
\end{ejercicio}

\begin{ejercicio}{2}{eje:3-cumple}
Hay $n$ personas, además de usted, con cumpleaños uniformes e independientes entre 365 días. (a) Escriba la probabilidad de que alguna comparta cumpleaños con usted y halle el menor $n$ para el que supera $1/2$. (b) Explique por qué ese $n$ es mucho mayor que 23 y por qué es mayor que $365/2$.
\fuente{adaptado de MIT 18.05~\cite{mit1805}, \emph{Problem Set}~1, problema~3, apartado~(c); el apartado~(b) es nuevo.}
\end{ejercicio}

\begin{ejercicio}{2}{eje:3-urna}
Una urna contiene $r$ bolas rojas y $v$ verdes, $r,v\ge1$. Se extraen dos sin reemplazamiento. (a) Demuestre, con un espacio muestral de resultados ordenados, que la probabilidad de que la segunda sea verde es igual a la de que lo sea la primera. (b) Si $r+v=16$ y es igual de probable que las dos bolas sean del mismo color que de colores distintos, halle $r$ y $v$.
\fuente{adaptado de Harvard Stat 110~\cite{stat110}, \emph{Homework}~1, problema~6.}
\end{ejercicio}

\begin{ejercicio}{2}{eje:3-geometrica}
Se eligen dos números $x$ e $y$ al azar en $[0,1]$, de modo que la probabilidad de cada región del cuadrado $[0,1]^{2}$ es su área. Calcule $\PP(\abs{x-y}>1/3)$ y $\PP(x+y<1/2)$.
\end{ejercicio}

\begin{ejercicio}{3}{eje:3-bonferroni}
Demuestre que, para sucesos cualesquiera, $\PP(A\cap B)\ge\PP(A)+\PP(B)-1$ y, más en general, $\PP(A\cap B\cap C)\ge\PP(A)+\PP(B)+\PP(C)-2$.
\fuente{adaptado de MIT 6.041~\cite{mit6041f10}, \emph{Problem Set}~1, problema~6.}
\end{ejercicio}

\begin{ejercicio}{3}{eje:3-interseccion}
Demuestre que, para sucesos $A_1,\dots,A_n$, $\PP\bigl(\bigcap_{i=1}^{n}A_i\bigr)\ge1-\sum_{i=1}^{n}\PP(A_i^{c})$. Aplíquelo: si cada uno de 20 componentes falla con probabilidad $0{,}001$, sin ninguna hipótesis de independencia, ¿qué puede asegurarse de la probabilidad de que no falle ninguno?
\end{ejercicio}

\begin{ejercicio}{3}{eje:3-distancia}
Sea $d(A,B)=\PP(A\dsim B)$, donde $A\dsim B=(A\setminus B)\cup(B\setminus A)$. (a) Demuestre que $d(A,B)=\PP(A)+\PP(B)-2\PP(A\cap B)$. (b) Demuestre que $d(A,C)\le d(A,B)+d(B,C)$.
\end{ejercicio}

\begin{ejercicio}{3}{eje:3-hash}
Se asignan $n$ claves a $N$ casillas de una tabla \emph{hash}, de forma uniforme e independiente. (a) Usando la \cref{prop:t3-cumple}, demuestre que si $n(n-1)\ge 2N\ln 2$, la probabilidad de alguna colisión es al menos $1/2$. (b) Para $N=2^{32}$, halle el menor $n$ que cumple esa condición y compárelo con $\sqrt{N}$.
\end{ejercicio}

\begin{ejercicio}{3}{eje:3-laplace}
(a) Demuestre que el modelo de Laplace (\cref{def:t3-laplace}) cumple los tres axiomas de Kolmogorov. (b) Demuestre que no existe ninguna ley de probabilidad sobre $\N$ que asigne la misma probabilidad a todos los sucesos elementales $\{n\}$.
\end{ejercicio}
